\section{Further research questions with exact targets}
\label{sec:research}

The following problems are organized around the structures proved above.
Each asks for an identity, a finite reduction, or a proof of a specified
conjecture.  A successful result should identify the precise function
class and normalization before a computation is interpreted.

\begin{enumerate}[leftmargin=*,label=\textbf{R\arabic*.}]
\item \textbf{Complete every fixed-order reconstruction family.}
Theorem~\ref{ert:main} supplies a finite decision procedure once $a$
is fixed. For $a=5$, the exact threshold $B_5=1920$ reduces the
smaller-order-ten problem to $6\le b<1920$; the adjacent case
$b=6$ and the gap-four case $b=7$ are already proved. Complete the
remaining cases using exact arithmetic, then apply the same procedure
at $a=6$, corresponding to smaller order twelve. A structural
factorization or a congruence may eliminate most of the finite range
before direct determinants are formed.  The target is a complete
all-$b$ theorem for each family, with every exceptional index accounted
for, rather than a larger table of experiments.

\item \textbf{Remove the fixed-order qualification.}
The leading positive determinant gives eventual nonvanishing but does
not control the lower Bernoulli degrees uniformly at small separation.
Find a determinant identity, an arithmetic invariant, or a recurrence
which prevents cancellation for all $a<b$.  The existing adjacent
and gap-four cases, together with the complete small-order families,
provide exact consistency constraints.  A proof uniform in $a$ would
settle the remaining unrestricted optimal-prefix conjecture.

\item \textbf{Give closed coefficient rows for higher-order reconstruction.}
The ordinary identities \eqref{r8c:ordinary} use a completely specified
rational inverse matrix.  Seek a product formula or a fixed-size kernel
for its rows, comparable to the known smaller-order inverse construction.
It should specify every polynomial coefficient directly and retain all
subtraction constants.  A formal reference to an unevaluated inverse
or to unknown minors is a weaker endpoint than such a closed kernel.

\item \textbf{Find the fifth- and higher-power tail cancellation.}
For each integer $d\ge5$, continue
$\sum_{n\ge0}\psi'(n+1)^d x_n^{-s}$ on its regular parity lattice.
The quartic boundary identity \eqref{qtm:boundary} suggests combining
the parity components of successive powers before applying finite
summation.  Determine which lower tail powers survive, identify a fixed
coordinate space for each $d$, and construct its anchored generator.
The dimension and the required weights should be proved, not inferred
from the first few moments.

\item \textbf{Treat the resonant parity as a separate spectral family.}
For the quartic family, the points $s=-3-2k$ are actual poles with
explicit residues $h_k$.  Determine their Laurent constants and
higher spectral coefficients by a normalization compatible with
\eqref{qtm:ordinary}.  These coefficients will involve logarithmic
moments, so the regular even-moment theorem cannot simply be
differentiated at its discrete values.  An anchored analytic generator
for the Laurent family would be a natural next object.

\item \textbf{Determine the functional depth of the quartic generator.}
The present construction gives a finite depth-at-most-three expression.
Specify allowable transformations of arguments and products, then ask
whether the particular depth-three combination reduces further.
A symbolic relation must be checked by a functional derivative
identity and an anchored endpoint.  Formal symbol or shuffle rank
should not be interpreted as arithmetic irreducibility of a value.

\item \textbf{Extend exact reconstruction to mixed and shifted tails.}
Replace the identical fourth power by
$\prod_j\zeta(s_j,n+a_j)$, beginning with positive integer orders
and a common rational shift.  Determine which parity cancellations
survive and the finite replacement for
$\zeta(-2m,1/2)=0$.  Character twists should be handled only after
the shift-dependent counterterms are explicit.  The goal is ordinary
convergent identities with a specified finite arithmetic coordinate set.

\item \textbf{Reduce the spectral logarithmic quadrature on special loci.}
For rational $a$ and selected algebraic masses, seek finite reductions
of \eqref{msj:eq:quadrature} and \eqref{one:digammaintegral} into
polylogarithm order derivatives, Gamma derivatives, or a precisely
defined multiple-Gamma family.  The ordinary integral is already an
exact answer at arbitrary parameters.  A finite special-function
reduction is an additional theorem and should not be assumed from
the first-derivative factorization.

\item \textbf{Classify noninteger-weight resonance crossings.}
Replace a polynomial monomial by $x^\rho$ in the one-sided master.
As $\rho$ crosses an odd integer, an argument of the Hurwitz function
passes through one.  Construct a joint Laurent formula in $s$ and
$\rho$ which specializes to \eqref{one:resonantC} and distinguishes
different orders of taking limits.  The local coefficient at each
crossing must match the polynomial theorem exactly.

\item \textbf{Build primitive formulas at every spectral order.}
The primitive \eqref{one:primitive} treats the nonlinear second jet
with a completely fixed dilogarithmic anchor.  Higher logarithmic
powers lead naturally to multiple polylogarithms after the same
rational substitution.  Develop a finite integration rule for each
$m$ and combine it with the classical primitives of
$\gamma_{m-2}(a)$.  Also determine when differentiation in a mass
commutes with an anchored shift primitive without an additional
endpoint term.

\item \textbf{Project the Stieltjes identities to arithmetic characters.}
At $a=p/q$, apply finite character sums to the all-order correction
formulas.  Give an exact dictionary between the resulting combinations
of $\gamma_j(p/q)$, Hurwitz derivatives, and Dirichlet-$L$ derivatives.
The reflected even and odd channels should be retained separately.
The first useful deliverable would cover small denominators with
the full logarithmic normalization in $q$.

\item \textbf{Return to the Gaussian period identities with new functional input.}
The unresolved $S_6$ and revised $S_8$ identities require a relation
that survives the known restrictions of the explored Cayley relation
systems.  The higher-depth polylogarithms appearing in the quartic
generator offer additional functional expressions, but no connection
to the target period is established here.  A useful next result would
give an explicit map from one such functional identity to the exact
canonical residual, including all lower-weight and branch terms.

\item \textbf{Formalize the finite and analytic interfaces.}
The Newton triangularization, Bernoulli modular recurrence, finite
summation identities, and resonant coefficient extraction are suitable
finite proof-assistant targets.  The separate analytic obligations are
normal convergence of the subtracted families and their agreement
with the original series on a common half-plane.  A formalization
should state those interfaces explicitly.  Validated interval
evaluation of the remaining quadratures is a different project from
formal proof and can be added independently.
\end{enumerate}
